% !PRINTED 591
\weberpage{591}
%<<CONT
kommen noch die associirten Zahlen. Aus 5. folgt dann noch,
dass eine Primzahl $p$ nur auf \gesperrt{eine Art} in die Summe von zwei
Quadraten zerlegt werden kann. Durch diese Zerlegung sind
aber die Zahlen $a, b$ noch nicht völlig bestimmt, sondern sie
können noch mit einander vertauscht und mit zwei Vorzeichen
versehen werden. Das kommt darauf hinaus, dass man $\pi$ durch
jede der vier associirten Zahlen $a+bi, -a-bi, -b+ai,$
$b-ai$ ersetzen kann. Ist $p$ ungerade, so muss von den beiden
Zahlen $a, b$ die eine gerade, die andere ungerade sein. Wählen
wir etwa für $b$ die gerade der beiden Zahlen, und bestimmen
das Vorzeichen so, dass $a \equiv 1 \pmod{4}$ wird, so ist unter den
vier associirten Zahlen eine bestimmte ausgewählt, die man die
\gesperrt{primäre nennen kann}\footnote{Gauss giebt an der erwähnten Stelle zwei verschiedene Bestimmungen
 für die primären Zahlen zur Auswahl, von denen dies die erste
ist. Er behält weiterhin die zweite bei.}.

Diese Definition der primären Zahlen lässt sich auf alle
ganzen Zahlen des Körpers $R(i)$ übertragen, deren Norm ungerade
 ist, und man hat dann das Gesetz, dass das Product
zweier primären Zahlen wieder eine primäre Zahl ergiebt.

Das System von vier associirten Zahlen im Körper $R(i)$ ist
analog dem Paar entgegengesetzter Zahlen im Körper $R$. In $R$
betrachtet man die positiven Zahlen als die primären.

Durch die Formel (27) des vorigen Paragraphen ist für
irgend eine gegebene Primzahl $p$ einer der Factoren $a+bi$ aus
der Kreistheilung abgeleitet. Nach den Formeln \S.\,172, (38), (39)
ist diese Zahl primär, wenn $p \equiv 5 \pmod{8}$, dagegen der primären
 entgegengesetzt, wenn $p \equiv 1 \pmod{8}$ ist. Darüber aber,
welche von den beiden conjugirten Zahlen $a \pm bi$ durch diese
Formeln dargestellt ist, haben wir kein allgemeines Kennzeichen.

Beispielsweise sind die complexen Primzahlen in $R(i)$, deren
Normen kleiner als 200 sind:
\[
\begin{aligned}
&1+i,\ 1+2i,\ 3+2i,\ 1+4i,\ 5+2i,\ 1+6i,\ 5+4i,\\
&7+2i,\ 5+6i,\ 3+8i,\ 5+8i,\ 9+4i,\ 1+10i,\ 3+10i,\\
&7+8i,\ 11+4i,\ 7+10i,\ 11+6i,\ 13+2i,\ 9+10i,\ 7+12i,\\
&1+14i.
\end{aligned}
\]
